% Shared style for (MATH101, 102, 311, 312) Calculus and Analysis.
% Notation and mathematical writing follow the current (MATH000) Mathematical Language.

\usepackage[a4paper,margin=1in,headheight=14pt,headsep=18pt]{geometry}
\usepackage[T1]{fontenc}
\usepackage{lmodern}
\usepackage{microtype}
\usepackage[intlimits]{amsmath}
\usepackage{amssymb,amsthm}
\usepackage{mathtools}
\usepackage{xcolor}
\usepackage{array}
\usepackage{booktabs}
\usepackage{longtable}
\usepackage{enumitem}
\usepackage{fancyhdr}
\usepackage{tocloft}
\usepackage{needspace}
\usepackage{etoolbox}
\usepackage{environ}
\usepackage[hidelinks]{hyperref}

% Institution identity.
\definecolor{POSTECHRed}{RGB}{166,25,85}
\newcommand{\institutiontext}[2]{\texorpdfstring{\textcolor{#1}{#2}}{#2}}
\DeclareRobustCommand{\POSTECH}{\institutiontext{POSTECHRed}{POSTECH}}
\DeclareRobustCommand{\PohangUniversityOfScienceAndTechnology}{%
  \institutiontext{POSTECHRed}{Pohang University of Science and Technology}%
}

% Canonical notation currently needed by the document shell.
% Add further notation here only after checking the current MATH000 standard.
\newcommand{\N}{\mathbb{N}}
\newcommand{\Nzero}{\mathbb{N}_0}
\newcommand{\Z}{\mathbb{Z}}
\newcommand{\Q}{\mathbb{Q}}
\newcommand{\R}{\mathbb{R}}
\newcommand{\C}{\mathbb{C}}
\newcommand{\dd}{\mathop{}\!\mathrm{d}}
\DeclareMathOperator{\dom}{dom}
\DeclareMathOperator{\codom}{codom}
\DeclareMathOperator{\id}{id}
\DeclareMathOperator{\im}{im}
\newcommand{\ptwto}{\to}
\newcommand{\unifto}{\rightrightarrows}
\newcommand{\ordinaryqed}{\ensuremath{\square}}
\newcommand{\majorqed}{\ensuremath{\blacksquare}}
\newcommand{\exceptionalqed}{\textsc{Q.E.D.}}

% Page typography.
\setlength{\parskip}{0pt}
\setlength{\parindent}{1.5em}
\linespread{1.0}
\setlist[itemize]{leftmargin=1.4em,itemsep=0.15em,topsep=0.25em}
\setlist[enumerate]{leftmargin=1.6em,itemsep=0.15em,topsep=0.25em}

% Tables reserved for Notation and Conventions as needed.
\newcolumntype{L}[1]{>{\raggedright\arraybackslash}p{#1}}
\setlength{\LTpre}{0.5em}
\setlength{\LTpost}{0.62em}
\newenvironment{notationtable}{%
  \par\vspace{0.05em}\small
  \renewcommand{\arraystretch}{1.18}%
  \begin{longtable}{@{}L{0.27\textwidth}L{0.65\textwidth}@{}}%
  \toprule
  Notation & Meaning \\
  \midrule
  \endfirsthead
  \toprule
  Notation & Meaning \\
  \midrule
  \endhead
  \bottomrule
  \endlastfoot
}{\end{longtable}}
\newenvironment{ruletable}{%
  \par\vspace{0.05em}\small
  \renewcommand{\arraystretch}{1.18}%
  \begin{longtable}{@{}L{0.24\textwidth}L{0.68\textwidth}@{}}%
  \toprule
  Item & Convention \\
  \midrule
  \endfirsthead
  \toprule
  Item & Convention \\
  \midrule
  \endhead
  \bottomrule
  \endlastfoot
}{\end{longtable}}

% Contents list only the top-level divisions.
\setlength{\cftbeforepartskip}{0.52em}
\renewcommand{\cftpartfont}{\normalsize}
\renewcommand{\cftpartpagefont}{\normalsize}

% Course-local numbering: 1.x, 2.x, 3.x, 4.x.
\newcounter{course}
\newcounter{mathitem}
\newcounter{problemitem}
\renewcommand{\themathitem}{\arabic{course}.\arabic{mathitem}}
\renewcommand{\theproblemitem}{\arabic{course}.\arabic{problemitem}}

\newcommand{\preliminarypart}[1]{%
  \clearpage
  \setcounter{course}{0}%
  \setcounter{mathitem}{0}%
  \setcounter{problemitem}{0}%
  \phantomsection
  \addcontentsline{toc}{part}{0\space #1}%
  \thispagestyle{plain}%
  \vspace*{0.48cm}%
  {\noindent\fontsize{20.5}{24.5}\selectfont\bfseries 0\quad #1\par}%
  \vspace{0.56cm}%
  \hrule height 0.4pt
  \vspace{0.78cm}%
}

\newcommand{\coursepart}[2]{%
  \clearpage
  \setcounter{course}{#1}%
  \setcounter{mathitem}{0}%
  \setcounter{problemitem}{0}%
  \phantomsection
  \addcontentsline{toc}{part}{#1\space #2}%
  \thispagestyle{plain}%
  \vspace*{0.48cm}%
  {\noindent\fontsize{20.5}{24.5}\selectfont\bfseries #1\quad #2\par}%
  \vspace{0.56cm}%
  \hrule height 0.4pt
  \vspace{0.78cm}%
}

% Optional course references. Empty arguments print nothing.
\newcommand{\relatedbooks}[1]{%
  \ifstrempty{#1}{}{%
    {\noindent\small Related Books: #1\par}%
    \vspace{0.62cm}%
  }%
}

% Mathematical environments share one course-local counter.
\newtheoremstyle{casplain}
  {6pt}{6pt}{\itshape}{}{\bfseries}{}{\newline}
  {\thmname{#1}\thmnumber{ #2}\thmnote{. #3}}
\newtheoremstyle{casdefinition}
  {6pt}{6pt}{}{}{\bfseries}{}{\newline}
  {\thmname{#1}\thmnumber{ #2}\thmnote{. #3}}
\newtheoremstyle{cassameline}
  {6pt}{6pt}{}{}{\itshape}{.}{0.5em}
  {\thmname{#1}\thmnumber{ #2}\thmnote{. #3}}

\theoremstyle{casdefinition}
\newtheorem{definition}[mathitem]{Definition}
\newtheorem{example}[mathitem]{Example}
\newtheorem{counterexample}[mathitem]{Counterexample}
\newtheorem{formula}[mathitem]{Formula}
\theoremstyle{casplain}
\newtheorem{theorem}[mathitem]{Theorem}
\newtheorem{proposition}[mathitem]{Proposition}
\newtheorem{lemma}[mathitem]{Lemma}
\newtheorem{corollary}[mathitem]{Corollary}
\theoremstyle{cassameline}
\newtheorem*{remark}{Remark}

\renewenvironment{proof}[1][Proof]
  {\par\Needspace{3\baselineskip}\noindent\textit{#1.}\hspace{0.5em}\ignorespaces}
  {\hfill$\square$\par}

\newcommand{\casblockspace}{\Needspace{4\baselineskip}}
\BeforeBeginEnvironment{definition}{\casblockspace}
\BeforeBeginEnvironment{example}{\casblockspace}
\BeforeBeginEnvironment{counterexample}{\casblockspace}
\BeforeBeginEnvironment{formula}{\casblockspace}
\BeforeBeginEnvironment{theorem}{\casblockspace}
\BeforeBeginEnvironment{proposition}{\casblockspace}
\BeforeBeginEnvironment{lemma}{\casblockspace}
\BeforeBeginEnvironment{corollary}{\casblockspace}
\BeforeBeginEnvironment{remark}{\casblockspace}

% Problems and Solutions are rendered separately from one paired source.
\newif\ifCASolutions
\CASolutionsfalse
\NewEnviron{problem}[1][]{%
  \refstepcounter{problemitem}%
  \ifCASolutions
  \else
    \Needspace{4\baselineskip}%
    \par\vspace{0.55em}%
    \noindent\textbf{Problem \theproblemitem}%
    \ifstrempty{#1}{}{\textbf{. #1}}%
    \par\vspace{0.28em}%
    \BODY
    \par\vspace{0.58em}%
  \fi
}
\NewEnviron{solution}{%
  \ifCASolutions
    \Needspace{4\baselineskip}%
    \par\vspace{0.55em}%
    \noindent\textit{Solution \theproblemitem.}\hspace{0.5em}%
    \BODY\hfill$\square$
    \par\vspace{0.58em}%
  \fi
}

% Lightweight opening shared by all three PDFs.
\newcommand{\CASOpening}[1]{%
  \pagenumbering{gobble}%
  \thispagestyle{empty}%
  \vspace*{0.72cm}%
  {\noindent\fontsize{18.5}{22}\selectfont\bfseries
    (MATH101, 102, 311, 312) Calculus and Analysis\par}%
  \vspace{0.12cm}%
  {\noindent\fontsize{11.8}{14.4}\selectfont\itshape #1\par}%
  \vspace{0.44cm}%
  {\noindent\normalsize Jungwoo Kim\par}%
  \vspace{0.02cm}%
  {\noindent\normalsize Department of Mathematics, \PohangUniversityOfScienceAndTechnology\par}%
  \vspace{0.44cm}%
  \hrule height 0.4pt
  \vspace{0.36cm}%
  {\noindent\large\bfseries Contents\par}%
  \vspace{0.18cm}%
  \begingroup
    \makeatletter
    \renewcommand{\contentsname}{}%
    \vspace{-2.1em}%
    \tableofcontents
    \makeatother
  \endgroup
  \thispagestyle{empty}%
  \clearpage
  \pagenumbering{arabic}%
}

\newcommand{\CASDocumentType}{Calculus and Analysis}
\pagestyle{fancy}
\fancyhf{}
\fancyhead[L]{\small\textit{Calculus and Analysis}}
\fancyhead[R]{\small\textit{\CASDocumentType}}
\fancyfoot[C]{\small\thepage}
\renewcommand{\headrulewidth}{0.35pt}
\renewcommand{\footrulewidth}{0pt}
\fancypagestyle{plain}{%
  \fancyhf{}
  \fancyhead[L]{\small\textit{Calculus and Analysis}}
  \fancyhead[R]{\small\textit{\CASDocumentType}}
  \fancyfoot[C]{\small\thepage}
  \renewcommand{\headrulewidth}{0.35pt}
  \renewcommand{\footrulewidth}{0pt}%
}
